%%%%%%%%%%%%%%%%%%%% solutions.tex %%%%%%%%%%%%%%%%%%%%%%%%%%%%%
% Solutions to Checkpoint Assessment Questions & Chapter Exercises
%%%%%%%%%%%%%%%%%%%%%%%%%%%%%%%%%%%%%%%%%%%%%%%%%%%%%%%%%%%%%

\chapter*{Solutions to Checkpoint Questions \& Exercises}
\addcontentsline{toc}{chapter}{Solutions to Checkpoint Questions \& Exercises}
\label{ch:solutions}

\section*{Chapter 1: The Multi-Core Communication Crisis \& Interconnect Physics}

\begin{description}
    \item[Question 1.1 (Tier 1: Fundamental)] \textbf{Answer: B}.
    According to the Elmore delay formulation (Eq.~\ref{eq:elmore_delay}), the propagation delay of a shared multi-drop wire is:
    \begin{equation*}
    \tau_{\mathrm{Elmore}} \approx 0.38 R_{\mathrm{wire}} C_{\mathrm{wire}} + R_{\mathrm{driver}} \sum_{i=1}^N C_{\mathrm{tap}, i} + \sum_{i=1}^N R_{\mathrm{seg}, i} C_{\mathrm{tap}, i}
    \end{equation*}
    Because each attached core introduces transceiver parasitic capacitance ($C_{\mathrm{taps}} \propto N$) and the physical wire length across the chip perimeter scales as $L \propto \sqrt{N}$ (increasing wire resistance $R_{\mathrm{wire}} \propto \sqrt{N}$ and wire capacitance $C_{\mathrm{wire}} \propto \sqrt{N}$), the cross-term product surges as $O(N^{3/2})$ to $O(N^2)$. For $N > 8$, wire charging time exceeds the processor clock period, resulting in catastrophic operating frequency throttling.

    \item[Question 1.2 (Tier 2: Architectural Analysis)] \textbf{Answer: B}.
    The Mayadas-Shatzkes effect (Eq.~\ref{eq:mayadas_shatzkes}) models electron scattering at polycrystalline grain boundaries in nanoscale copper lines. When trench width approaches the electron mean free path ($\lambda_0 \approx 39.0\,\mathrm{nm}$ in bulk copper), grain boundary reflection ($R \approx 0.3$--$0.5$) and surface diffuse reflection combine with non-conductive diffusion barrier liners (Ta/TaN) to surge effective resistivity by $5\times$ to $15\times$ above bulk values, destroying global interconnect RC performance.

    \item[Question 1.3 (Tier 3: Design Trade-Off)] \textbf{Answer: A}.
    An $N \times N$ crossbar requires $N^2$ crosspoint multiplexer switches and $2N \times W_{\mathrm{data}}$ wiring tracks crossing the central layout. For $N = 128$ and $W_{\mathrm{data}} = 64$ bits, this dictates 16,384 crosspoints and 16,384 parallel global wire tracks, creating insurmountable physical routing track congestion, crosstalk jitter, and quadratic silicon area explosion.
\end{description}

\section*{Chapter 2: Switching Techniques, Flow Control \& Buffer Management}

\begin{description}
    \item[Question 2.1 (Tier 1: Fundamental)] \textbf{Answer: B}.
    Wormhole switching operates at the granularity of individual flits rather than entire packets. By allowing body flits to stream across established router paths behind the head flit, input buffer depth can be reduced to as few as 1--4 flits, reducing router SRAM buffer area by 75\%--93\% compared to Store-and-Forward and Virtual Cut-Through.

    \item[Question 2.2 (Tier 2: Architectural Analysis)] \textbf{Answer: B}.
    The credit turnaround loop Round-Trip Time is:
    \begin{equation*}
    \mathrm{RTT} = t_{\mathrm{tx}} + t_{\mathrm{flight\_fwd}} + t_{\mathrm{rx}} + t_{\mathrm{flight\_cred}} = 1 + 1 + 1 + 1 = 4\,\mathrm{cycles}
    \end{equation*}
    By the Fundamental Throughput Theorem (Eq.~\ref{eq:bmin_theorem}), sustaining 100\% wire-rate throughput requires buffer depth $B \ge \mathrm{RTT} = 4$ flits. With $B = 2$, the transmitter exhausts its credits after 2 cycles and is forced to stall during cycles 3 and 4 while waiting for credit feedback, throttling link utilization to $50\%$.

    \item[Question 2.3 (Tier 3: Design Trade-Off)] \textbf{Answer: B}.
    In the SWIFT NoC, lightweight tokens travel ahead of data flits to reserve crossbar switch allocation in advance. When a token is granted, incoming data flits cut directly through the crossbar into the outgoing link without writing into or reading from input FIFO SRAM cells, saving up to 38\% of dynamic router power at saturation.
\end{description}

\section*{Chapter 3: High-Performance Router Microarchitecture \& Arbitration}

\begin{description}
    \item[Question 3.1] \textbf{Answer: B}.
    The canonical synchronous 5-stage router pipeline comprises: Buffer Write (BW), Route Computation (RC), Virtual Channel Allocation (VA), Switch Allocation (SA), and Switch Traversal (ST), followed by inter-router Link Traversal (LT).

    \item[Question 3.2] \textbf{Answer: A}.
    When Client $k$ receives a grant ($G_k = 1$), row $k$ is cleared to zero ($M[k][j] = 0\ \forall j \ne k$) and column $k$ is set to one ($M[i][k] = 1\ \forall i \ne k$). This places client $k$ at the lowest priority relative to all other requesters, implementing an exact, single-cycle Least Recently Used (LRU) policy.

    \item[Question 3.3] \textbf{Answer: B}.
    Updating round-robin pointers in later iterations would cause output grant arbiters to synchronize into identical priority positions when traffic is heavy. Restricting pointer updates to the first iteration preserves pointer diversity (entropy), preventing cyclic grant lockouts and achieving 100\% throughput under uniform traffic.
\end{description}

\section*{Chapter 4: Direct, Indirect \& High-Radix Interconnect Topologies}

\begin{description}
    \item[Bisection Bandwidth Bound] Under uniform random traffic, destinations are chosen independently and uniformly with probability $1/N$. Exactly half ($N/2$) of the destinations lie across any bisection cut that partitions the $N$ nodes into two equal halves. Thus, exactly half ($50\%$) of all injected packets must cross the bisection. Each node injects at rate $\Theta$, yielding $(N/2)\Theta$ packets per half; half cross the cut: $T_{\mathrm{cut}} = \frac{N \Theta}{4}$. Because total bidirectional channel capacity across the cut is $B_C$, stability requires $\frac{N \Theta}{4} \le \frac{B_C}{2}$, which proves the fundamental upper bound:
    \begin{equation*}
    \Theta \le \frac{2 B_C}{N}
    \end{equation*}

    \item[Folded Torus Delay] In a standard $8 \times 8$ torus, neighbor links have length $L_0 = 1.2\,\mathrm{mm}$, while the wrap-around link spans $(k-1) L_0 = 7 \times 1.2\,\mathrm{mm} = 8.4\,\mathrm{mm}$. Under the Elmore quadratic delay model ($\tau \propto L^2$), the wrap-around delay is:
    \begin{equation*}
    \frac{\tau_{\mathrm{wrap}}}{\tau_{\mathrm{local}}} = \left(\frac{8.4\,\mathrm{mm}}{1.2\,\mathrm{mm}}\right)^2 = 7^2 = 49\times
    \end{equation*}
    In a folded torus, interleaved node placement ($0, 7, 1, 6, 2, 5, 3, 4$) equalizes all link lengths to $2 L_0 = 2.4\,\mathrm{mm}$. The longest link delay is $(2.4/1.2)^2 = 4\times$, delivering a $12.25\times$ delay reduction ($49/4$) and eliminating timing closure violations.

    \item[Clos Non-Blocking Sizing] For a 3-stage Clos network $C(n, m, r)$ with $N = 64$ and $n = 8$ (yielding $r = N/n = 8$ input switches):
    \begin{itemize}
        \item Strictly non-blocking: by Clos's theorem (1953), $m \ge 2n - 1 = 2(8) - 1 = 15$ middle switches.
        \item Rearrangeably non-blocking: by the Slepian-Duguid theorem, $m \ge n = 8$ middle switches.
    \end{itemize}

    \item[High-Radix Trade-Offs] In low-radix 2D meshes ($k = 4$), the average hop count scales as $O(\sqrt{N})$. Every intermediate router hop expends dynamic energy in FIFO buffer writes/reads, crossbar traversal, and link charging. In a high-radix Dragonfly or Flattened Butterfly network ($k \ge 32$), the network diameter is compressed to $D \le 3$. Even though each individual router crossbar has higher radix, packet energy is dominated by hop count: traversing 2 high-radix routers consumes far less total silicon energy than traversing 12--16 intermediate low-radix routers.
\end{description}

\section*{Chapter 5: Deadlock Theory, Turn Models \& Adaptive Routing}

\begin{description}
    \item[CDG Construction] A $2 \times 2$ mesh contains 8 directed channels: $\{c_{01}, c_{10}, c_{02}, c_{20}, c_{13}, c_{31}, c_{23}, c_{32}\}$. Minimal fully adaptive routing permits clockwise turns ($c_{01} \to c_{13} \to c_{32} \to c_{20} \to c_{01}$) and counter-clockwise turns ($c_{02} \to c_{23} \to c_{31} \to c_{10} \to c_{02}$), generating directed cycles in the Channel Dependency Graph (CDG). West-First routing prohibits all turns ending in the West direction ($N \to W$ and $S \to W$). This breaks both circular wait cycles, transforming the CDG into a Directed Acyclic Graph (DAG) and guaranteeing deadlock freedom.

    \item[Duato Theorem Dynamics] Duato's theorem mandates that the escape subnetwork $R_1$ ($VC_0$) forms an acyclic connected sub-graph. If flits could freely transition from the escape subnetwork $VC_0$ back into adaptive channels $VC_1$, an adaptive cycle could propagate back into $VC_0$, creating an indirect circular wait ($VC_0 \to VC_1 \to VC_1 \to VC_0$). Prohibiting transitions from $VC_0$ back to $VC_1$ ensures that $VC_0$ acts as a guaranteed drain sink for any congested packet.

    \item[Protocol Deadlock Scenario] In a MESI directory protocol, a core injects a Read Request ($Req$). The home directory core must issue Invalidation Requests ($Inv$) to multiple sharer cores and collect Invalidation Acks ($Ack$) before injecting the Data Response ($Resp$). If $Req$ and $Resp$ share the same physical router buffers, heavy $Req$ traffic fills all input buffers. The directory cannot send $Resp$ because buffers are saturated; cores cannot accept new $Req$ because their transaction queues are awaiting $Resp$. Although XY routing prevents physical routing deadlocks, circular wait develops across message classes ($Req \to Inv \to Ack \to Resp$). Dedicating separate Virtual Networks (VNETs) to each class breaks this protocol deadlock.
\end{description}

\section*{Chapter 6: Low-Power Architectures \& Dark Silicon Management}

\begin{description}
    \item[Breakeven Time Calculation]
    \begin{itemize}
        \item Re-activation energy overhead: $E_{\mathrm{overhead}} = C_{\mathrm{virtual}} \cdot V_{dd}^2 = 3.2\,\mathrm{pF} \times (0.85\,\mathrm{V})^2 = 2.312\,\mathrm{pJ}$.
        \item Static leakage power: $P_{\mathrm{leak}} = V_{dd} \cdot I_{\mathrm{leak}} = 0.85\,\mathrm{V} \times 8.5\,\mathrm{mA} = 7.225\,\mathrm{mW}$.
        \item Breakeven time: $\tau_{\mathrm{breakeven}} = E_{\mathrm{overhead}} / P_{\mathrm{leak}} = \frac{2.312\,\mathrm{pJ}}{7.225\,\mathrm{mW}} = 0.320\,\mathrm{ns}$.
        \item At $1.5\,\mathrm{GHz}$ ($T_{\mathrm{clk}} = 0.667\,\mathrm{ns}$), $\tau / T_{\mathrm{clk}} = 0.320 / 0.667 \approx 0.48$ cycles. Sleeping for just **1 idle clock cycle** saves net energy.
    \end{itemize}

    \item[DVFS Power Scaling] Dynamic power scales as $P_{\mathrm{dyn}} \propto V_{dd}^2 \cdot f$.
    Baseline: $P_1 \propto 1.0^2 \times 1.5 = 1.50$.
    Scaled: $P_2 \propto 0.75^2 \times 1.0 = 0.5625 \times 1.0 = 0.5625$.
    Percentage reduction: $\frac{1.50 - 0.5625}{1.50} = \mathbf{62.5\%}$ dynamic power reduction.

    \item[Lookahead Wakeup Hazards] If a flit enters a buffer whose virtual rail has only charged to $50\%$ ($0.425\,\mathrm{V}$), CMOS gate propagation delays surge exponentially as $(V_{gs} - V_{th}) \to 0$. Transistors enter the near-threshold regime, causing severe setup-time violations, metastable latch states, and corrupted flit payloads in SRAM cells and crossbar muxes.
\end{description}

\section*{Chapter 7: Thermal Modeling, Hotspots \& Fault Tolerance}

\begin{description}
    \item[Thermal Time Constant]
    \begin{itemize}
        \item $\tau_{\mathrm{th}} = R_v \cdot C_{\mathrm{th}} = 15.0\,\mathrm{K/W} \times 3.6 \times 10^{-3}\,\mathrm{J/K} = 54.0\,\mathrm{ms}$.
        \item At $2.5\,\mathrm{GHz}$ ($T_{\mathrm{clk}} = 0.4\,\mathrm{ns}$), the number of clock cycles is $N_{\mathrm{cyc}} = \frac{54.0 \times 10^{-3}}{0.4 \times 10^{-9}} = 135 \times 10^6$ cycles.
        \item Physical explanation: silicon and copper bulk materials possess large thermal heat capacities, making thermal diffusion 5 orders of magnitude slower than microarchitectural packet bursts ($10^2$--$10^3$ cycles). Reactive throttling only engages long after hotspot temperature limits have been exceeded.
    \end{itemize}

    \item[Electromigration \& Black's Equation]
    \begin{itemize}
        \item $\frac{E_a}{k_B} = \frac{0.90 \times 1.6022 \times 10^{-19}}{1.3806 \times 10^{-23}} \approx 10,443\,\mathrm{K}$.
        \item $\frac{1}{T_2} - \frac{1}{T_1} = \frac{1}{358.15} - \frac{1}{333.15} = -0.0002097\,\mathrm{K^{-1}}$.
        \item $\frac{\mathrm{MTTF}(T_2)}{\mathrm{MTTF}(T_1)} = \exp(10,443 \times -0.0002097) = e^{-2.190} = 0.1119$ ($\mathbf{11.2\%}$).
        \item Percentage reduction: $100\% - 11.2\% = \mathbf{88.8\%}$ reduction in interconnect lifetime from a $25^\circ\mathrm{C}$ temperature rise.
    \end{itemize}

    \item[Non-Gaussian Sensor Filtering] High-frequency power supply noise ($L \cdot di/dt$ ground bounce) creates heavy-tailed impulse noise outliers ($>5^\circ\mathrm{C}$). Standard linear Kalman filters minimize quadratic error, causing large outliers to skew the gain matrix and diverge. Huber thresholding replaces quadratic loss with linear penalty for errors exceeding threshold $\delta$, bounding outlier influence and stabilizing state estimation.

    \item[ZoneDefense Deadlock Freedom] Concave fault patterns force packets into detour loops that reverse dimension routing order, introducing $Y \to X$ turns in a 2D mesh. Grouping arbitrary faults into convex bounding rectangles guarantees that detour perimeters remain strictly monotonic and continuous, preserving acyclic CDG properties without virtual channels.
\end{description}

\section*{Chapter 8: Three-Dimensional Networks-on-Chip \& TSVs}

\begin{description}
    \item[TSV Elmore Delay vs. Planar Wire]
    \begin{itemize}
        \item TSV: $\tau_{\mathrm{tsv}} = R_{\mathrm{driver}} C_{\mathrm{tsv}} + R_{\mathrm{tsv}} C_{\mathrm{tsv}} / 2 = 400 \times 120.1 \times 10^{-15} + 0.0541 \times 120.1 \times 10^{-15} / 2 \approx \mathbf{48.04\,\mathrm{ps}}$.
        \item Planar wire: $\tau_{\mathrm{planar}} = R_{\mathrm{driver}} C_{\mathrm{wire}} + R_{\mathrm{wire}} C_{\mathrm{wire}} / 2 = 400 \times 270 \times 10^{-15} + 120 \times 270 \times 10^{-15} / 2 = 108\,\mathrm{ps} + 16.2\,\mathrm{ps} = \mathbf{124.2\,\mathrm{ps}}$.
        \item Speedup: $\frac{\tau_{\mathrm{planar}}}{\tau_{\mathrm{tsv}}} = \frac{124.2\,\mathrm{ps}}{48.04\,\mathrm{ps}} = \mathbf{2.59\times}$ faster propagation delay.
    \end{itemize}

    \item[3D-IC Thermal Stacking Calculation]
    \begin{itemize}
        \item Heat sink dissipation: $T_0 = T_{\mathrm{amb}} + R_{\mathrm{hs}} (P_0 + P_1 + P_2) = 40.0^\circ\mathrm{C} + 0.40\,\mathrm{K/W} \times 75.0\,\mathrm{W} = \mathbf{70.0^\circ\mathrm{C}}$.
        \item Middle tier: $T_1 = T_0 + R_{\mathrm{inter}} (P_1 + P_2) = 70.0^\circ\mathrm{C} + 0.35\,\mathrm{K/W} \times 50.0\,\mathrm{W} = \mathbf{87.5^\circ\mathrm{C}}$.
        \item Top tier: $T_2 = T_1 + R_{\mathrm{inter}} P_2 = 87.5^\circ\mathrm{C} + 0.35\,\mathrm{K/W} \times 25.0\,\mathrm{W} = \mathbf{96.25^\circ\mathrm{C}}$.
        \item Top tier is $96.25 - 70.0 = \mathbf{26.25^\circ\mathrm{C}}$ hotter due to the cumulative inter-tier bonding resistance ladder.
    \end{itemize}

    \item[TSV Keep-Out Zone (KOZ)] Copper has a thermal expansion coefficient ($\alpha_{\mathrm{Cu}} \approx 17 \times 10^{-6}\,\mathrm{K^{-1}}$) $6\times$ higher than silicon ($\alpha_{\mathrm{Si}} \approx 2.6 \times 10^{-6}\,\mathrm{K^{-1}}$). During post-bonding cooling, thermal contraction mismatch induces high compressive hoop stress ($>200\,\mathrm{MPa}$) in silicon. By the piezoresistive effect, this alters silicon electron and hole effective masses, shifting threshold voltage ($V_{th}$) and mobility by over $15\%$. Transistors must be placed outside the $3$--$5\,\mu\mathrm{m}$ KOZ to prevent timing failures.

    \item[Vertical Crossbar Scaling] A 5-port 2D router requires $5 \times 5 = 25$ crosspoints, while a 7-port 3D router requires $7 \times 7 = 49$ crosspoints ($+96\%$ area and quadratic wire loading). Vertical arbitration-free fast-paths bypass two-stage switch allocation by assigning dedicated pipeline registers to Up/Down channels, eliminating vertical arbitration contention.
\end{description}

\section*{Chapter 9: Multi-Chiplet Interconnects \& UCIe Standard}

\begin{description}
    \item[The Giant Pizza Yield Problem]
    \begin{itemize}
        \item Monolithic $800\,\mathrm{mm^2}$ ($A = 8.0\,\mathrm{cm^2}$, $A \cdot D_0 = 0.64$):
        $Y = \left(\frac{1 - e^{-0.64}}{0.64}\right)^2 = \left(\frac{1 - 0.5273}{0.64}\right)^2 = \mathbf{54.55\%}$.
        \item Four independent $200\,\mathrm{mm^2}$ chiplets ($A = 2.0\,\mathrm{cm^2}$, $A \cdot D_0 = 0.16$):
        $Y_{\mathrm{chiplet}} = \left(\frac{1 - e^{-0.16}}{0.16}\right)^2 = \left(\frac{1 - 0.8521}{0.16}\right)^2 = \mathbf{85.45\%}$.
        \item Physical explanation: A single dust defect on a wafer destroys an entire expensive $800\,\mathrm{mm^2}$ die. With smaller $200\,\mathrm{mm^2}$ chiplets, that defect only scraps one small piece while the other three chiplets survive with $85.5\%$ yield, slashing net manufacturing cost by $40\%$--$60\%$.
    \end{itemize}

    \item[Buffer Sizing via Little's Law]
    \begin{itemize}
        \item $\mathrm{RTT}_{\mathrm{credit}} = 3.2\,\mathrm{ns} + 0.4\,\mathrm{ns} + 3.2\,\mathrm{ns} = 6.8\,\mathrm{ns}$. In clock cycles ($T_{\mathrm{clk}} = 0.4\,\mathrm{ns}$): $\mathrm{RTT}_{\mathrm{credit}} = 6.8 / 0.4 = \mathbf{17\,\mathrm{cycles}}$.
        \item By Little's Law: $B_{\mathrm{min}} \ge \lambda \cdot \mathrm{RTT} = 1 \times 17 = \mathbf{17\,\mathrm{flits}}$.
        \item With only 8 flits of buffer, the transmitter exhausts credits after 8 cycles and stalls for 9 cycles, causing link utilization to collapse to $\frac{8}{17} = \mathbf{47.1\%}$.
    \end{itemize}

    \item[The Four-Way Traffic Jam] When four packets block each other in a circular wait ($N \to E \to S \to W \to N$), diverting one packet into a temporary side register (the In/Outbound swap buffer) breaks the circular wait condition. The remaining three packets traverse the intersection unblocked; the diverted packet then rejoins the link without dropping any data flits.

    \item[Energy Comparison] For $1\,\mathrm{TB}$ ($8 \times 10^{12}\,\mathrm{bits}$):
    Organic substrate ($1.20\,\mathrm{pJ/bit}$): $E_{\mathrm{org}} = 8 \times 10^{12} \times 1.20 \times 10^{-12} = \mathbf{9.60\,\mathrm{J}}$.
    Silicon interposer ($0.33\,\mathrm{pJ/bit}$): $E_{\mathrm{inter}} = 8 \times 10^{12} \times 0.33 \times 10^{-12} = \mathbf{2.64\,\mathrm{J}}$.
    Net energy saved: $9.60 - 2.64 = \mathbf{6.96\,\mathrm{J}}$ ($72.5\%$ reduction).
\end{description}

\section*{Chapter 10: Optical \& Wireless Networks-on-Chip}

\begin{description}
    \item[The Optical Tuning Fork]
    \begin{itemize}
        \item Optical round-trip length: $L = 2 \pi R = 2 \pi \times 6.0\,\mu\mathrm{m} = \mathbf{37.699\,\mu\mathrm{m}}$.
        \item Resonant mode order: $m = \frac{L \cdot n_{\mathrm{eff}}}{\lambda_0} = \frac{37,699 \times 2.45}{1550} = \mathbf{60}$.
        \item Physical explanation: A $1540\,\mathrm{nm}$ wave does not fit an exact integer number of cycles inside the cavity; it cannot constructively interfere and bypasses the ring entirely along the through waveguide.
    \end{itemize}

    \item[The Laser Flashlight]
    \begin{itemize}
        \item Insertion loss: $L_{\mathrm{path}} = 3.5 \times 0.25 + 16 \times 0.005 + 30 \times 0.12 + 20 \times 0.015 + 2 \times 0.50 + 1.50 + 3.00 = \mathbf{10.355\,\mathrm{dB}}$.
        \item Required optical power: $P_{\mathrm{laser}} = -21.0\,\mathrm{dBm} + 10.355\,\mathrm{dB} = -10.645\,\mathrm{dBm} = \mathbf{0.0862\,\mathrm{mW}}$ ($86.2\,\mu\mathrm{W}$).
        \item Electrical wall-plug power ($\eta_{\mathrm{wpe}} = 18\%$): $P_{\mathrm{elec}} = \frac{0.0862}{0.18} = \mathbf{0.479\,\mathrm{mW}}$.
    \end{itemize}

    \item[The Hot CPU]
    \begin{itemize}
        \item Wavelength drift: $\Delta \lambda = 20^\circ\mathrm{C} \times 0.10\,\mathrm{nm/^\circ C} = \mathbf{2.0\,\mathrm{nm}}$.
        \item Heating power per MRR: $2.0\,\mathrm{nm} \times 20.0\,\mathrm{mW/nm} = 40.0\,\mathrm{mW}$.
        \item Total tuning power for 128 MRRs: $128 \times 40.0\,\mathrm{mW} = \mathbf{5.12\,\mathrm{W}}$.
    \end{itemize}

    \item[Walkie-Talkie vs. Whisper] Unicast electrical packet replication requires sending 63 sequential packets through intermediate routers, incurring dozens of pipeline stages and crossbar arbitrations ($>20$--$50$ cycles). A $60\,\mathrm{GHz}$ wireless antenna radiates electromagnetic waves across the die at the speed of light in silicon dioxide ($c / \sqrt{\epsilon_r} \approx 150\,\mu\mathrm{m/ps}$), broadcasting the invalidation message to all 63 cores simultaneously in a **single clock cycle** ($<1\,\mathrm{ns}$).
\end{description}

\section*{Chapter 11: Hardware Security \& Trust in On-Chip Networks}

\begin{description}
    \item[The Thermal Trojan \& Junction Melting]
    \begin{itemize}
        \item Under idle traffic:
        \begin{equation*}
        T_{\mathrm{idle}} = T_{\mathrm{amb}} + P_{\mathrm{base}} \cdot R_{\mathrm{th}} = 45^\circ\mathrm{C} + 1.6 \times 12.5 = 45 + 20 = \mathbf{65.0^\circ\mathrm{C}}
        \end{equation*}
        \item Under Trojan attack ($P_{\mathrm{total}} = 1.6 + 4.8 = 6.4\,\mathrm{W}$):
        \begin{equation*}
        T_{\mathrm{attack}} = T_{\mathrm{amb}} + P_{\mathrm{total}} \cdot R_{\mathrm{th}} = 45^\circ\mathrm{C} + 6.4 \times 12.5 = 45 + 80 = \mathbf{125.0^\circ\mathrm{C}}
        \end{equation*}
        \item The thermal RC time constant is $\tau_{\mathrm{th}} = R_{\mathrm{th}} C_{\mathrm{th}} = 12.5 \times 2.0 \times 10^{-3} = 0.025\,\mathrm{s} = 25.0\,\mathrm{ms}$. \\
        The temperature evolution follows $T(t) = T_{\mathrm{attack}} - (T_{\mathrm{attack}} - T_{\mathrm{idle}}) e^{-t / \tau_{\mathrm{th}}}$. Setting $T(t) = 95^\circ\mathrm{C}$:
        \begin{equation*}
        95 = 125 - (125 - 65) e^{-t / 25} \implies 60 e^{-t / 25} = 30 \implies e^{-t / 25} = 0.5
        \end{equation*}
        Solving for $t$: $t_{\mathrm{trigger}} = 25 \times \ln(2) = 25 \times 0.6931 = \mathbf{17.33\,\mathrm{ms}}$. Emergency DTM throttles within $17.3\,\mathrm{ms}$.
    \end{itemize}

    \item[Zero-Multiplier Shift-and-Add Arithmetic]
    \begin{itemize}
        \item With $\alpha = 2^{-3} = 1/8$, $\mu_t = \mu_{t-1} + \frac{1}{8}(T_t - \mu_{t-1})$. The division by 8 is synthesized as an arithmetic right shift by 3 bits (\texttt{>>> 3}), replacing an area-heavy multiplier with 0 logic gates (hardwired bit shifts).
        \item $3 \cdot D_{t-1} = (D_{t-1} \ll 1) + D_{t-1}$, which requires exactly one wire routing shift and one adder.
        \item To preserve 3 fractional bits for $\alpha = 1/8$ across an 8-bit temperature input ($0\text{--}255$), the fixed-point format requires at least $8 + 3 = 11$ bits (format \texttt{UQ8.3}). Adding sign margin for difference error requires a 12-bit signed accumulator to avoid overflow.
    \end{itemize}

    \item[Timing Covert Channel Capacity]
    \begin{itemize}
        \item $\mathrm{SNR} = \frac{\Delta \tau^2}{\sigma_{\mathrm{jitter}}^2} = \frac{12^2}{6^2} = \frac{144}{36} = \mathbf{4.0}$. \\
        The covert capacity upper bound is $C_{\mathrm{covert}} = \frac{1}{2} \log_2(1 + 4.0) = \frac{1}{2} \log_2(5) \approx \mathbf{1.161\,\mathrm{bits/packet}}$.
        \item With LFSR jitter injected, total noise variance is $\sigma_{\mathrm{total}}^2 = 6^2 + 36^2 = 36 + 1296 = 1332$. \\
        $\mathrm{SNR}_{\mathrm{masked}} = \frac{144}{1332} \approx 0.1081$. \\
        New capacity: $C_{\mathrm{masked}} = \frac{1}{2} \log_2(1 + 0.1081) \approx \mathbf{0.074\,\mathrm{bits/packet}}$.
        \item Bandwidth reduction: $\frac{1.161 - 0.074}{1.161} = \mathbf{93.6\%}$ reduction.
    \end{itemize}

    \item[Interconnect Intrinsic PUFs]
    \begin{itemize}
        \item Because the two crossbar paths are symmetrically drawn, their delay difference $\Delta t = t_{\mathrm{path\_A}} - t_{\mathrm{path\_B}}$ is distributed normally as $\mathcal{N}(0, 2\sigma^2) = \mathcal{N}(0, 32\,\mathrm{ps}^2)$ with mean zero. By symmetry of the Gaussian distribution around zero, $P(\Delta t < 0) = \mathbf{0.50}$ (50.0\%).
        \item A 50\% probability of producing bit 0 or 1 yields maximal Shannon entropy: $H(R) = - (0.5 \log_2 0.5 + 0.5 \log_2 0.5) = \mathbf{1.0\,\mathrm{bit/cell}}$.
        \item The variations stem from atomic-scale line edge roughness and random dopant fluctuations (quantum-scale fabrication jitter) that are physically unique to each die and impossible to reproduce or reverse-engineer from the GDSII layout.
    \end{itemize}
\end{description}

\section*{Chapter 12: Domain-Specific NoCs for LLMs, Transformers \& SNNs}

\begin{description}
    \item[Ring All-Reduce vs. In-Network Reduction Trees]
    \begin{itemize}
        \item In a Ring All-Reduce with $N = 64$ nodes and $M = 256\,\mathrm{MB} = 2.048 \times 10^9\,\mathrm{bits}$: \\
        Total data transferred per node: $V_{\mathrm{node}} = 2 \frac{63}{64} \times 2.048 \times 10^9 = 4.032 \times 10^9\,\mathrm{bits}$. \\
        Link transmission time at $100\,\mathrm{Gb/s}$:
        \begin{equation*}
        t_{\mathrm{ring}} = \frac{4.032 \times 10^9\,\mathrm{bits}}{100 \times 10^9\,\mathrm{bits/s}} = \mathbf{0.04032\,\mathrm{s}} = \mathbf{40.32\,\mathrm{ms}}
        \end{equation*}
        \item In a 4-ary In-Network Reduction Tree (height $H = 3$), nodes transmit local slices ($M/N = 4\,\mathrm{MB} = 3.2 \times 10^7\,\mathrm{bits}$) up the tree where routers sum them in-flight, and the root broadcasts the reduced $4\,\mathrm{MB}$ slices down. Total node transmission volume is $2 \times (M/N) = 6.4 \times 10^7\,\mathrm{bits}$.
        \begin{equation*}
        t_{\mathrm{tree}} = \frac{6.4 \times 10^7}{100 \times 10^9} + H \cdot t_{\mathrm{hop\_latency}} \approx 0.00064\,\mathrm{s} + 3 \times 1.5\,\mathrm{ns} \approx \mathbf{0.640\,\mathrm{ms}}
        \end{equation*}
        \item Speedup: $S = \frac{40.32\,\mathrm{ms}}{0.640\,\mathrm{ms}} \approx \mathbf{63.0\times}$ across the tensor slice!
    \end{itemize}

    \item[Ruche Express Link Wire Physics]
    \begin{itemize}
        \item Resistance per millimeter ($L = 1000\,\mu\mathrm{m}$): \\
        For M3 ($w = 0.050\,\mu\mathrm{m}$): number of squares is $1000 / 0.050 = 20{,}000\,\square$. \\
        $R_{\mathrm{wire, M3}} = 20{,}000 \times 0.35\,\Omega = \mathbf{7{,}000\,\Omega/mm}$. \\
        For M8 ($w = 0.500\,\mu\mathrm{m}$): number of squares is $1000 / 0.500 = 2{,}000\,\square$. \\
        $R_{\mathrm{wire, M8}} = 2{,}000 \times 0.035\,\Omega = \mathbf{70\,\Omega/mm}$ ($100\times$ lower!).
        \item For a $3.0\,\mathrm{mm}$ link ($3000\,\mu\mathrm{m}$): \\
        $R_{\mathrm{M3}} = 21{,}000\,\Omega$, $C_{\mathrm{M3}} = 3000 \times 0.22\,\mathrm{fF} = 660\,\mathrm{fF}$. \\
        $\tau_{\mathrm{Elmore, M3}} = 0.38 \times 21{,}000 \times 660 \times 10^{-15} = 5.266 \times 10^{-9}\,\mathrm{s} = \mathbf{5{,}266\,\mathrm{ps}}$ ($5.27\,\mathrm{ns}$). \\
        $R_{\mathrm{M8}} = 210\,\Omega$, $C_{\mathrm{M8}} = 3000 \times 0.15\,\mathrm{fF} = 450\,\mathrm{fF}$. \\
        $\tau_{\mathrm{Elmore, M8}} = 0.38 \times 210 \times 450 \times 10^{-15} = 3.591 \times 10^{-11}\,\mathrm{s} = \mathbf{35.91\,\mathrm{ps}}$.
        \item Clock period is $T_{\mathrm{clk}} = 500\,\mathrm{ps}$. The M8 Ruche wire traverses $3\,\mathrm{mm}$ (skipping 4 router tiles) in $35.9\,\mathrm{ps} \ll 500\,\mathrm{ps}$ within a single clock cycle, whereas the M3 wire requires more than 10 clock cycles without repeaters!
    \end{itemize}

    \item[Hybrid Bonding \& The LLM Memory Wall]
    \begin{itemize}
        \item Weights to read per token: $70 \times 10^9 \times 2\,\mathrm{bytes} = 140\,\mathrm{GB/token}$. \\
        Required memory bandwidth: $140\,\mathrm{GB/token} \times 80\,\mathrm{tokens/s} = \mathbf{11{,}200\,\mathrm{GB/s}} = \mathbf{11.2\,\mathrm{TB/s}}$.
        \item Standard HBM3 provides $3.2\,\mathrm{TB/s} < 11.2\,\mathrm{TB/s}$. An HBM3 system cannot sustain $80\,\mathrm{tokens/s}$ (it is throttled to at most $3.2 / 0.14 = \mathbf{22.8\,\mathrm{tokens/s}}$).
        \item The hybrid-bonded system delivers $16.4\,\mathrm{TB/s}$. Supported token generation throughput is:
        \begin{equation*}
        \mathrm{Throughput}_{\max} = \frac{16.4\,\mathrm{TB/s}}{0.14\,\mathrm{TB/token}} = \mathbf{117.1\,\mathrm{tokens/second}}
        \end{equation*}
    \end{itemize}

    \item[AER Spike Multicasting \& SNNcut Hypergraph Partitioning]
    \begin{itemize}
        \item Firing count: $1{,}024 \times 0.10 = 102.4 \approx 102$ neurons. \\
        Under naive unicast routing, each firing neuron injects 256 individual packets: \\
        $\mathrm{Total}_{\mathrm{unicast}} = 102 \times 256 = \mathbf{26{,}112\,\mathrm{packets}}$.
        \item Under SNNcut hypergraph partitioning, target neurons are grouped such that each axon targets only 6 cores on average, transmitting a single multicast packet per target core: \\
        $\mathrm{Total}_{\mathrm{SNNcut}} = 102 \times 6 = \mathbf{612\,\mathrm{packets}}$.
        \item Bandwidth reduction: $\frac{26{,}112 - 612}{26{,}112} = \frac{25{,}500}{26{,}112} = \mathbf{97.65\%}$ reduction in network packet volume!
    \end{itemize}
\end{description}
